\printedpage{19}
% head: A. Grothendieck : Dérivateurs, Ch. XV, Édition des Universités de Montpellier II et Paris VII
% foot: Version 5 novembre 2010

\msp{page 28}
\textsc{Proof}. Everything has been proved, if we recall that the objects of $\mathcal{M}^r$ are $r^{\pi_0}$-accessible (\textit{cf.} (3.2)).

\section*{4\quad Complements on Quillen structures}

\textbf{Definition 4.1.} Let $\mathcal{M}$ be a category. By a \textit{Quillen precouple} in $\mathcal{M}$ we mean a pair $(\Phi,\Psi)$ of subsets of $\Fl(\mathcal{M})$, satisfying the following two conditions.
\begin{itemize}[leftmargin=2.5em]
\item[\textit{a})] We have $\Phi\longleftrightarrow\Psi$, \textit{i.e.} $\Psi\subset\Phi_*$, or what amounts to the same, $\Phi\subset\Psi^*$ -- or again: for every commutative square in $\mathcal{M}$
\[
\begin{tikzcd}[row sep=large, column sep=large]
A \arrow[r, "u"] \arrow[d, "i"'] & X \arrow[d, "p"] \\
B \arrow[ur, dotted, "w" description] \arrow[r, "v"] & Y \;,
\end{tikzcd}
\tag*{(4.1.1)}
\]
with $i\in\Phi$, $p\in\Psi$, there exists $w$ such that $wi=u$, $pw=v$.
\item[\textit{b})] Every $f\in\Fl(\mathcal{M})$ factors as
\[
f \;=\; p\,i\;, \qquad \text{with } i\in\Phi,\, p\in\Psi\ .
\]
\end{itemize}
We write
\[
\Phi \;\qpre\; \Psi
\tag*{(4.1.2)}
\] % ?? uncertain: the printed relation sign -- a double arrow with a small vertical bar in the middle (rendered by \qpre)
to indicate that $(\Phi,\Psi)$ is a Quillen precouple.

We say that the Quillen precouple is \textit{closed}, if $\Phi$ is left closed, $\Psi$ right closed. We shall see below (cor. 2, page 31) that this implies (hence is equivalent to) % [sic] clos à droit -> clos à droite, gender agreement with « à droite »
\[
\Phi \;=\; \Psi^*\;, \qquad \Psi \;=\; \Phi_*\ .
\tag*{(4.1.3)}
\]

\msp{page 29}
A closed Quillen precouple is also called a Quillen couple.

\textbf{Proposition 4.2.} \textit{Let} $(\Phi,\Psi)$ \textit{be a Quillen precouple in} $\mathcal{M}$\textit{, and let} $j\in\Fl(\mathcal{M})$\footnote[9]{\textbf{N.B.} We use only that $\Phi\longleftrightarrow\Psi$, and that $j$ can be put in the form $pi$, with $i\in\Phi$, $p\in\Psi$.}. % footnote mark printed as "(9)"
\textit{We then have the equivalent conditions on} $j$ \textit{:}
\begin{itemize}[leftmargin=2.5em]
\item[a)] $j\in\widetilde{\Phi}$.
\item[b)] $j\in\Psi^*$.
\item[c)] $j$ \textit{is a direct factor of an} $i\in\Phi$.
\item[c$'$)] $j$ \textit{is a direct factor of an} $i\in\Phi$ \textit{having the same source as} $j$\textit{, more precisely, there exists a diagram}
\[
\begin{tikzcd}[row sep=large, column sep=large]
A \arrow[r, "j"] \arrow[d, shift right=1ex, "\id_A"'] & B \arrow[d, shift right=1ex, "q"'] \\
A \arrow[u, dotted, shift right=1ex, "\id_A"'] \arrow[r, "i"] & B' \;, \arrow[u, dotted, shift right=1ex, "p"']
\end{tikzcd}
\tag*{(4.2.1)}
\]
\printedpage{20}
% head: A. Grothendieck : Dérivateurs, Ch. XV, Édition des Universités de Montpellier II et Paris VII
% foot: Version 5 novembre 2010
\textit{with} $i\in\Phi$\textit{, and}
\[
qj \;=\; i\;, \qquad pi \;=\; j\;, \qquad pq \;=\; \id_B\;,
\tag*{(4.2.2)}
\]
\textit{so that} $j$ \textit{is a direct factor of} $i\in\Phi$\textit{, via} $(\id_A,q)$\textit{, admitting the retraction} $(\id_A,p)$.
\item[d)] \textit{For every commutative solid-arrow diagram}
\[
\begin{tikzcd}[column sep=small]
A \arrow[rr, "j"] \arrow[dr, "i"'] & & B \arrow[dl, dotted, bend left=40, "q"] \\
& B' \arrow[ur, "p"'] &
\end{tikzcd}
\;,
\tag*{(4.2.3)}
\]
i.e. \textit{every factorization of} $j$ \textit{as} $pi$\textit{, if we have} $p\in\Psi$\textit{, there exists a section} $q$ \textit{of} $B'$ \textit{over} $B$ \textit{in} $A\backslash\mathcal{M}$\textit{,} i.e. \textit{an arrow} $q:B\longrightarrow B'$ \textit{such that}
\[
pq \;=\; \id_B \qquad \textit{and} \qquad qj \;=\; i\ .
\]
\end{itemize}

\textsc{Proof} of the proposition. It proceeds by the following sequence of implications, all $(\pm)$ tautological:
\[
(c') \;\Longrightarrow\; (c) \;\Longrightarrow\; (a) \;\Longrightarrow\; (b) \;\Longrightarrow\; (d) \;\Longrightarrow\; (c')\;,
\]
\msp{page 30}
of which the first four use only the hypothesis $\Phi\longleftrightarrow\Psi$ (from $(a)$ of definition 4.1), and only the last, $(d)\Longrightarrow(c')$, calls on the Quillen factorization hypothesis ($(b)$ in 4.1); and even then it suffices to assume it for $j$ only!

$(c')\Longrightarrow(c)$ pure tautology, $(c)\Longrightarrow(a)$ since we know that $\widetilde{\Phi}$ is stable under direct factors (and contains $\Phi$), $(a)\Longrightarrow(b)$ since $\widetilde{\Phi}\subset\Psi^*$ ($\widetilde{\Phi}$ being by definition the smallest left closed subset of $\Fl(\mathcal{M})$ containing $\Phi$, and $\Psi^*$ being left closed), $(b)\Longrightarrow(d)$ since we write the given diagram (4.2.2) (solid arrows) as % [sic] (4.2.2) -> (4.2.3), the diagram given in condition d) is (4.2.3)
\[
\begin{tikzcd}[row sep=large, column sep=large]
A \arrow[r, "u"] \arrow[d, "j\in\Psi^*"'] & X \arrow[d, "p\in\Psi"] \\
B \arrow[ur, dotted, "q" description] \arrow[r, "v"] & Y \;,
\end{tikzcd}
\tag*{$(*)$}
\]
putting $Y=B$, $v=\id_B$ in it, and rebaptizing $X$ as $B'$, and $u$ as $i$: the hypothesis $(b)$, $j\in\Psi^*$, means that in a solid-arrow diagram $(*)$, there always exists $q$ such that $qj=u$, $pq=v$, and expressing this in the particular case where $v=\id_B$, we find $(d)$.

It remains to prove $(d)\Longrightarrow(c')$, this time using the factorization hypothesis on $j$, \textit{i.e.} that we can write $j$ in the form $pi$:
\[
\begin{tikzcd}[column sep=small]
A \arrow[rr, "j"] \arrow[dr, "i"'] & & B \arrow[dl, dotted, bend left=40, "q"] \\
& B' \arrow[ur, "p"'] &
\end{tikzcd}
\;,
\]
with $i\in\Phi$, $p\in\Psi$. By virtue of $(d)$, since $p\in\Psi$, there thus exists $q:B\longrightarrow B'$ satisfying
\printedpage{21}
% head: A. Grothendieck : Dérivateurs, Ch. XV, Édition des Universités de Montpellier II et Paris VII
% foot: Version 5 novembre 2010
\msp{page 31}
\[
pq \;=\; \id_B\;, \qquad qj \;=\; i\;,
\]
which proves $(c')$, since $i\in\Phi$.

\textbf{Corollary 1.} \textit{Let} $(\Phi,\Psi)$ \textit{be a Quillen precouple. Then we have}
\[
\begin{aligned}
\widetilde{\Phi} \;\;=\;\; \Psi^* \;\;=\;\; &\textit{set of objects of } \underline{\Fl}(\mathcal{M}) \textit{ which are direct fac-}\\
&\textit{tors of elements of } \Phi\ ,
\end{aligned}
\tag*{(4.3)}
\] % printed line break of the display falls inside "fac-teurs"; kept as printed
\textit{and dually}
\[
\begin{aligned}
\overline{\Psi} \;\;=\;\; \Phi_* \;\;=\;\; &\textit{set of objects of } \underline{\Fl}(\mathcal{M}) \textit{ which are direct fac-}\\
&\textit{tors of elements of } \Psi\ .
\end{aligned}
\tag*{(4.4)}
\] % printed line break of the display falls inside "fac-teurs"; kept as printed

\textbf{Corollary 2.} \textit{The following conditions on the Quillen precouple} $(\Phi,\Psi)$ \textit{are equivalent.}
\begin{itemize}[leftmargin=2.5em]
\item[a)] \textit{The precouple is} closed, i.e. $\Phi=\widetilde{\Phi}$\textit{,} $\Psi=\overline{\Psi}$\textit{.}
\item[b)] $\Phi=\Psi^*$\textit{,} $\Psi=\Phi_*$\textit{.}
\item[c)] $\Phi$ \textit{and} $\Psi$ \textit{are stable under direct factors.}
\end{itemize}
Now let
\[
W \;\subset\; \Fl(\mathcal{M})
\tag*{(4.5)}
\]
be a localizer, satisfying the usual conditions
\[
\left\{
\begin{aligned}
&\textit{a})\quad W \text{ contains the isomorphisms.}\\[1ex]
&\textit{b})\quad \text{If } X\xrightarrow{\;f\;}Y\xrightarrow{\;g\;}Z \text{ in } \mathcal{M}\text{, if two of } f,\ g,\ gf \text{ are}\\
&\phantom{\textit{b})}\quad \text{in } W\text{, then the third too.}
\end{aligned}
\right.
\tag*{(4.6)}
\]
($W$ mildly saturated). Suppose moreover given four other subsets
\[
C\ ,\ TC\ , \qquad F\ ,\ TF
\tag*{(4.7)}
\]
of $\Fl(\mathcal{M})$, satisfying the following conditions.
\[
\left\{
\begin{aligned}
&\textit{a})\quad (C,TF) \text{ and } (TC,F) \text{ are Quillen precouples.}\\[1ex]
&\textit{b})\quad TC\subset C\cap W,\ TF\subset F\cap W.
\end{aligned}
\right.
\tag*{(4.8)}
\]

\msp{page 32}
These properties are spelled out as follows:
\begin{itemize}[leftmargin=2.5em]
\item[\textbf{M\,1}] Lifting-extension properties, $C\longleftrightarrow TF$ and $TC\longleftrightarrow F$.
\item[\textbf{M\,2}] Factorization property for every arrow $f$ of $\mathcal{M}$ as $f=pi$, with, at our choice, $i\in C$, $p\in TF$, or $i\in TC$, $p\in F$.
\item[\textbf{M$'$}] $TC\subset C\cap W$, $TF\subset F\cap W$.
\end{itemize}

\printedpage{22}
% head: A. Grothendieck : Dérivateurs, Ch. XV, Édition des Universités de Montpellier II et Paris VII
% foot: Version 5 novembre 2010
This is therefore part of the axioms of a weak Quillen structure for $W$ (XIII, p. 27, def. 1) -- the stability axioms M\,3, M\,4, and the ``properness'' (propriety) axioms M$''$ are missing, namely $F\subset\mathrm{Fib}_W$, $C\subset\mathrm{Cof}_W$.

\textbf{Proposition 4.3.} \textit{Let} $\mathcal{M}$ \textit{be a category,} $(C,TF)$ \textit{and} $(TC,F)$ \textit{two Quillen precouples in} $\mathcal{M}$\textit{, satisfying}
\[
TC \;\subset\; C\;, \qquad TF \;\subset\; F\;,
\tag*{(4.3.1)}
\]
\textit{and let} $W$ \textit{be a subset of} $\Fl(\mathcal{M})$ \textit{of which for the moment we assume only} 4.6, (b)\footnote[10]{We use only that if $X\xrightarrow{\;f\;}Y\xrightarrow{\;g\;}Z$, and if $gf\in W$, then $f\in W\Longleftrightarrow g\in W$. We do not need $W$ to contain the isomorphisms, nor to be stable under composition.}\textit{, and} % footnote mark printed as "(10)"
\[
W \;\supset\; TC\ ,\ TF\ ,
\tag*{(4.3.2)}
\]
\textit{so we have}
\[
TC \;\subset\; C\cap W\ , \qquad TF \;\subset\; F\cap W\ .
\tag*{(4.3.3)}
\]
\textit{Let us consider the following extension-lifting conditions} (a), (a$'$)\textit{, and the factorization conditions} (b), (b$'$), (c), (d) (\textit{recalling that} $\widetilde{TC}=F^*$\textit{,} $\overline{TF}=C_*$\textit{, by virtue of proposition} 4.2 (\textit{cor.} 1) \textit{applied to the Quillen precouples} $(TC,F)$\textit{,} $(C,TF)$):
\begin{itemize}[leftmargin=2.5em]
\item[a)] $C\cap W\subset\widetilde{TC}$\quad (\textit{hence} $TC\subset C\cap W\subset\widetilde{TC}$).
\item[a$'$)] $F\cap W\subset\overline{TF}$\quad (\textit{hence} $TF\subset F\cap W\subset\overline{TF}$).
\item[b)] $W\subset TF\circ\widetilde{TC}$.
\item[b$'$)] $W\subset\overline{TF}\circ TC$.
\item[c)] $W\subset\overline{TF}\circ\widetilde{TC}$.
\item[d)] $W\subset TF{\circ}TC$ (\textit{hence, by} (4.3.2)\textit{,} $W=TF{\circ}TC$ \textit{if} $W$ \textit{is stable under composition}).\footnote[11]{\textbf{N.B.} Condition $(a)$ also means that $(C\cap W,F)$ is still a Quillen precouple, resp. $(b)$ means that $(C,F\cap W)$ is one -- and then this pair of Quillen precouples this time satisfies $(d)$, \textit{i.e.} $W=TF'\circ TC'$.} % footnote mark printed as "(11)"
\end{itemize}

\msp{page 33}
\textit{Then all the conditions except the last are equivalent} (\textit{and they are therefore implied by the latter}). \textit{Thus when} $TC$ \textit{or} $TF$ \textit{are stable under direct factors} (\textit{e.g. if} $TC$ \textit{is left closed, or} $TF$ \textit{right closed}), \textit{all the conditions under consideration are equivalent} (\textit{since by corollary} 2 \textit{to proposition} 4.2\textit{, we then have} $TC=\widetilde{TC}$ \textit{resp.} $TF=\overline{TF}$\textit{, hence} (b$'$) \textit{resp.} (b) \textit{is then equivalent to} (d)).

\textsc{Proof}. It proceeds by the diagram of implications
\[
\begin{tikzcd}[column sep=small, row sep=small]
& (d) \arrow[dl, Rightarrow] \arrow[dr, Rightarrow] & \\
(b) \arrow[dr, Rightarrow] & \mathrm{Q} & (b') \arrow[dl, Rightarrow] \\
& (c) \arrow[dl, Rightarrow] \arrow[dr, Rightarrow] & \\
(a) \arrow[uu, Rightarrow] & & (a') \arrow[uu, Rightarrow]
\end{tikzcd}
\;,
\]
\printedpage{23}
% head: A. Grothendieck : Dérivateurs, Ch. XV, Édition des Universités de Montpellier II et Paris VII
% foot: Version 5 novembre 2010
where the four implications of the upper square Q are tautological. It thus suffices to establish the implications
\[
(c) \;\Longrightarrow\; (a) \;\Longrightarrow\; (b)\;,
\]
the implications $(c)\Longrightarrow(a')\Longrightarrow(b')$ follow from these by duality.

Suppose then $(c)$, \textit{i.e.}
\[
W \;\subset\; \overline{TF}\circ\widetilde{TC}\;,
\tag*{$(*)$}
\]
and let us prove $(a)$, \textit{i.e.} that
\[
C\cap W \;\subset\; \widetilde{TC}\ .
\]
But if $j\in C\cap W$, to prove that $j\in\widetilde{TC}$, let us write $j=pi$, $i\in\widetilde{TC}$, $p\in\overline{TF}$, which is possible by the hypothesis $(*)$ and $j\in W$.

\msp{page 34}
\[
\begin{tikzcd}[column sep=small]
A \arrow[rr, "j\in C"] \arrow[dr, "i\in\widetilde{TC}"'] & & B \arrow[dl, dotted, bend left=60, "q"] \\
& B' \arrow[ur, "p\in\overline{TF}"'] &
\end{tikzcd}
\;.
\]
Let us now use the hypothesis $j\in C$ and $C\longleftrightarrow\overline{TF}$, we thus find (prop. 3.2, $(a)\Longrightarrow(d)$) an arrow $q:B\longrightarrow B'$ such that $pq=\id_B$, $qj=i$, so $j$ appears as a direct factor of $i$. Since $i\in\widetilde{TC}$ and $\widetilde{TC}$ is stable under direct factors, we conclude $j\in\widetilde{TC}$, \quad \textit{q.e.d.}

$(a)\Longrightarrow(b)$. We thus assume
\[
C\cap W \;\subset\; \widetilde{TC}\;,
\tag*{$(*)$}
\]
and we want to prove $(b)$, \textit{i.e.}
\[
W \;\subset\; TF\circ\widetilde{TC}\ .
\]
Let then $f\in W$, and let us factor $f$ as $pi$, with $i\in C$, $p\in TF$. (O.K. because of the Quillen precouple $(C,TF)$). Since $TF\subset W$, we have $p\in W$, and since $pi\in W$, we conclude $i\in W$ (mild saturation of $W$ -- we do have to use it after all!). We thus have $i\in C{\cap}W$, hence $i\in\widetilde{TC}$ by $(*)$, whence $f\in TF\circ\widetilde{TC}$, \quad \textit{q.e.d.}

\textbf{Corollary 1.} \textit{Suppose} $TC$ \textit{stable under direct factors. Then all the conditions of proposition} 4.3 \textit{are equivalent, and are equivalent to the relation}
\[
C\cap W \;=\; TC\ .
\tag*{(4.3.4)}
\]
\textit{Dually, if} $TF$ \textit{is stable under direct factors, the conditions of proposition} 4.3 \textit{are all equivalent and are equivalent to}
\[
F\cap W \;=\; TF\ .
\tag*{(4.3.4$'$)}
\]